\chapter{O painel com efeitos fixos de dois sentidos}
\label{cap:painel}

Este é o principal instrumento de inferência do trabalho (decisão D8), e o que
mais provavelmente será cobrado em detalhe.

\section{A equação}

\begin{equation}
\Delta \mathrm{lulc}_{it} = \alpha_i + \gamma_t + \beta \, \Delta x_{it} + \varepsilon_{it}
\label{eq:2fe}
\end{equation}

onde:
\begin{itemize}
\item $\Delta \mathrm{lulc}_{it}$ é a \emph{variação anual} do uso da terra na
unidade $i$ no ano $t$;
\item $\alpha_i$ é o \textbf{efeito fixo de unidade};
\item $\gamma_t$ é o \textbf{efeito fixo de ano};
\item $\Delta x_{it}$ é a variação da explicativa de interesse;
\item $\varepsilon_{it}$ é o erro, com erros-padrão agrupados por unidade.
\end{itemize}

\section{O que cada efeito fixo faz, em português}

\textbf{$\alpha_i$, o efeito de unidade}, absorve tudo o que \emph{não muda no
tempo} naquele lugar: solo, relevo, distância a mercados, tradição produtiva.
E o ponto que dá poder ao estimador é que ele absorve isso \textbf{mesmo que
tais atributos não estejam observados no painel}. Você não precisa medir a
qualidade do solo de cada município para controlá-la; basta que ela seja
constante no tempo.

\textbf{$\gamma_t$, o efeito de ano}, absorve os choques comuns a todas as
unidades naquele ano: preços internacionais, câmbio, crises macroeconômicas.

Combinados, eles fazem o estimador \textbf{comparar cada unidade consigo mesma
ao longo do tempo}, líquida do que aconteceu com todo mundo naquele ano.
Tecnicamente, isso equivale a subtrair de cada observação a média da unidade e a
média do ano (a transformação \emph{within} de dois sentidos), e estimar sobre o
que sobra.

\section{O ponto fino: o que eles absorvem \emph{numa equação em diferenças}}

Aqui está a sutileza que separa quem entende de quem repete. Como as duas
pontas da equação entram em primeira diferença (D7), os efeitos fixos absorvem
coisa \textbf{distinta} da que absorveriam num modelo escrito em nível.

A diferenciação já retirou o \emph{patamar} de cada município. Então $\alpha_i$,
que num modelo em nível seria esse patamar, passa a retirar a
\textbf{tendência própria da unidade} --- o ritmo médio a que ela vinha mudando
ao longo da série. E $\gamma_t$ absorve o choque comum às \emph{variações}
daquele ano, e não ao estoque acumulado até ele.

\begin{quote}\itshape\color{cortinta}
O que resta para $\beta$ estimar é o quanto a variação de um município num ano
se afasta do que a tendência dele e o ano em curso já explicavam.
\end{quote}

\begin{nabanca}
``Como estou em primeira diferença, o efeito fixo de unidade não me tira o nível
--- a diferença já tirou. Ele me tira a tendência própria daquele município. O
que o $\beta$ mede é o desvio em relação ao que a tendência local e o ano em
curso já previam.''
\end{nabanca}

\section{O que o painel controla e o que ele não controla}

\begin{naodiz}
O painel com efeitos fixos \textbf{não} controla um determinante que
\emph{varie dentro da unidade ao longo do tempo} e que tenha ficado fora do
modelo. É exatamente por isso que os coeficientes deste trabalho são descritos
como \textbf{associações bem controladas}, e não como efeitos causais. \\[4pt]
Ele também não isola marcos específicos: o efeito fixo de ano absorve
\emph{todos} os choques comuns daquele ano, sem distingui-los entre si.
\end{naodiz}

\section{As três camadas de exigência}

Os instrumentos não respondem à mesma pergunta e obedecem a uma hierarquia
declarada. Tratá-los como equivalentes seria a origem mais provável de um
excesso de afirmação.

\vspace{6pt}
\begin{tabular}{@{}p{2.9cm}p{5.2cm}p{6.4cm}@{}}
\toprule
\textbf{Camada} & \textbf{O que responde} & \textbf{Limite} \\
\midrule
Correlação em nível estadual & Houve coincidência temporal entre as duas séries no estado? & De 11 a 39 observações e nenhum controle por confundidores \\
\addlinespace
Painel com efeitos fixos & Dentro de cada unidade, a variação de $X$ acompanha a de $Y$ depois de removidos os choques comuns do ano? & Não isola marco específico; não controla determinante que varie dentro da unidade no tempo \\
\addlinespace
Diferenças em diferenças & Este marco regulatório teve efeito diferencial em Goiás frente a estados de controle? & \textbf{Não identifica neste caso} \\
\bottomrule
\end{tabular}

\vspace{10pt}
\subsection{Por que a terceira camada foi rebaixada}

Este é um ponto que precisa ser dito com firmeza, porque é onde uma banca de
ciências ambientais mais provavelmente vai querer que você fale sobre política.

Os quatro marcos submetidos ao teste são a estabilização monetária do Plano
Real (tomada em 1995), o \emph{boom} de \emph{commodities} de 2003, o Código
Florestal de 2012 e a reorganização de mercado de 2018. \textbf{Nenhum deles é
goiano.} Dois são federais e dois são de alcance nacional ou internacional, e em
qualquer dos casos eles alcançaram no mesmo momento o estado tratado e os
controles do mesmo bioma.

Sem grupo não tratado, nem o par que sobrevive a todos os testes de robustez
--- tendências paralelas, estudo de evento e placebo temporal --- identifica
efeito causal.

\begin{quote}\itshape\color{cortinta}
O problema está no \emph{desenho}, não na estimação, e nenhum teste adicional o
corrige.
\end{quote}

\begin{nabanca}
``Rodei o desenho de diferenças em diferenças e ele passa em todos os testes de
robustez que se costuma exigir. Mesmo assim eu o rebaixei, porque a robustez não
resolve um problema de desenho: os marcos são federais e não deixam grupo não
tratado. O que aquele coeficiente mede é exposição diferencial ao mesmo choque,
e é assim que ele é reportado. Separar as hipóteses exigiria variação
\emph{dentro} do estado, entre municípios mais e menos expostos a cada marco, e
esse desenho eu não tenho.''
\end{nabanca}

\chapter{Erros-padrão: três réguas que não são intercambiáveis}
\label{cap:erros}

Um coeficiente sem erro-padrão não diz nada, e um erro-padrão calculado com a
régua errada diz coisa falsa com aparência de precisão. O trabalho usa três, e
declara qual está em uso em cada tabela.

\section{Agrupamento por unidade (\emph{cluster})}

Corrige a correlação \emph{dentro} de cada unidade ao longo do tempo. A
intuição: se um município tem um choque próprio que dura vários anos, suas 40
observações não são 40 informações independentes.

\section{Agrupamento nas duas dimensões}

Corrige também a correlação \emph{entre} unidades no mesmo ano. É a régua
adotada nas especificações que trazem efeito fixo de ano.

\begin{armadilha}[Um detalhe que a auditoria do próprio trabalho corrigiu]
Dispensar o efeito fixo de ano \textbf{não} obriga a abrir mão do agrupamento
por ano --- são escolhas independentes. Uma tabela do apêndice publicava um
bloco só na régua de unidade e descrevia isso como imposição do desenho; era
acoplamento de código. E recaía justamente onde mais custa, porque um regressor
constante dentro do ano tem o erro-padrão subestimado pelo agrupamento por
unidade. Corrigido, o coeficiente do estoque segue com $p<0{,}001$ e os três
sinais de demanda deixam de cruzar 5\%.
\end{armadilha}

\section{HAC de Newey-West, para séries agregadas}

Nas séries agregadas no nível estadual, em que o painel não se aplica, a
inferência usa erros robustos a heterocedasticidade \emph{e} autocorrelação, na
receita de Newey-West com duas defasagens. A matriz de covariância é

\begin{equation}
\hat{\Omega} = \hat{\Gamma}_0 + \sum_{\ell=1}^{L} w_\ell \left(\hat{\Gamma}_\ell + \hat{\Gamma}_\ell'\right),
\qquad w_\ell = 1 - \frac{\ell}{L+1},
\end{equation}

com $L = 2$: em vez de supor erros independentes, ela soma as autocovariâncias
das duas primeiras defasagens, com peso decrescente.

\begin{naodiz}
Com cerca de 38 observações anuais, o erro-padrão ingênuo subestima a incerteza
--- e o próprio teste corrigido tem \textbf{pouco poder} nesse tamanho de
amostra. Isso protege contra falso positivo e \emph{esconde efeitos reais
fracos}. As duas metades dessa frase precisam ser ditas juntas.
\end{naodiz}

\begin{nabanca}[A regra de ouro deste trabalho sobre réguas]
``Duas réguas de inferência não são comparáveis entre si. Onde eu tenho as duas,
reporto as duas lado a lado, e digo qual decide. Trocar de régua no meio de uma
conferência produziria uma robustez que é da régua, e não do dado.''
\end{nabanca}

\chapter{Estacionariedade, integração e regressão espúria}
\label{cap:integracao}

Este capítulo e o seguinte sustentam a decisão D16, que nasceu de um erro
próprio. Entendê-los bem é o que permite apresentar esse episódio como
credibilidade em vez de falha.

\section{O que é uma série integrada}

Uma série é \textbf{estacionária}, ou I(0), quando sua média e sua variância não
mudam ao longo do tempo e os choques se dissipam. Ela é \textbf{integrada de
ordem 1}, ou I(1), quando precisa ser diferenciada uma vez para ficar
estacionária --- o caso clássico é o passeio aleatório, $y_t = y_{t-1} +
\epsilon_t$, em que cada choque fica incorporado para sempre. Ela é I(2) quando
nem a primeira diferença é estacionária.

\section{Por que isso é fatal para uma regressão entre séries}

Duas séries I(1) independentes, geradas por processos que nada têm a ver um com
o outro, produzem regressões com $R^2$ alto e $t$ enorme com frequência muito
maior que os 5\% nominais. Isso é a \textbf{regressão espúria}: a estatística
não está medindo relação, está medindo o fato de que as duas séries vagam.

O mesmo vale para o teste de precedência: aplicado a séries integradas, ele
\textbf{fabrica precedência inexistente}.

\section{ADF contra KPSS: duas hipóteses nulas opostas}

Este é o ponto que a banca pode cobrar, e a resposta precisa ser exata.

\begin{itemize}
\item O \textbf{ADF} (Dickey-Fuller aumentado) tem por hipótese nula a
\emph{presença de raiz unitária}. Portanto $p$ \textbf{baixo} indica série
\textbf{estacionária}.
\item O \textbf{KPSS} tem por hipótese nula a \emph{estacionariedade}. Portanto
$p$ \textbf{baixo} indica o \textbf{contrário}.
\end{itemize}

Rodar os dois é confrontá-los. Quando concordam, o veredito é firme; quando
divergem, a divergência é registrada e não escondida.

\vspace{6pt}
\begin{tabular}{@{}lccl@{}}
\toprule
\textbf{Série} & \textbf{$p$ ADF} & \textbf{$p$ KPSS} & \textbf{Veredito (ADF)} \\
\midrule
Agricultura do Sul, em nível & $0{,}011$ & $0{,}010$ & estacionária \\
Pastagem do Norte, em nível & $0{,}960$ & $0{,}010$ & não estacionária \\
Agricultura do Sul, em 1\textsuperscript{a} dif. & $0{,}001$ & $0{,}050$ & estacionária \\
Pastagem do Norte, em 1\textsuperscript{a} dif. & $0{,}922$ & $0{,}010$ & não estacionária \\
Agricultura do Sul, em 2\textsuperscript{a} dif. & $<0{,}001$ & $0{,}100$ & estacionária \\
Pastagem do Norte, em 2\textsuperscript{a} dif. & $0{,}001$ & $0{,}100$ & estacionária \\
\bottomrule
\end{tabular}

\vspace{10pt}
\textbf{Como ler esta tabela em voz alta.} Os dois testes \emph{concordam} nas
duas linhas da pastagem do Norte --- ambos a dão não estacionária em nível e em
primeira diferença ---, e é dela que vem o diagnóstico. Eles \emph{divergem}
nas duas linhas da agricultura do Sul, que o ADF dá estacionária e o KPSS não.
E \emph{voltam a concordar} nas duas segundas diferenças. São estas as linhas
que fixam $d_{\max} = 2$: sem elas, a ordem de integração da pastagem ficaria
afirmada e não mostrada.

\begin{armadilha}[Não chame acordo de discordância]
Nas linhas da pastagem do Norte, ADF e KPSS \textbf{concordam}. É nas linhas da
agricultura do Sul que eles divergem. Trocar isso na fala inverte o sentido do
diagnóstico.
\end{armadilha}

\begin{naodiz}
O $p$ do KPSS é truncado na tabela de valores críticos da rotina: $0{,}010$ e
$0{,}100$ são \emph{limites}, e se leem como ``$\leq$'' e ``$\geq$''. \\[4pt]
E com 40 observações e duas quebras estimadas nessas mesmas séries, nenhum dos
dois testes é decisivo sozinho. A classificação vale como \emph{a resposta do
teste nesta amostra}, e não como propriedade estabelecida do processo. O que a
tabela estabelece com segurança é o \textbf{descasamento} entre as duas séries,
que é o que invalida o teste clássico.
\end{naodiz}

\section{A consequência lógica: séries de ordens diferentes não cointegram}

Se a agricultura do Sul é I(0) e a pastagem do Norte é I(2), não existe relação
de longo prazo entre os níveis para um teste detectar. Qualquer ``precedência''
encontrada nessa configuração é \textbf{artefato}, por construção.

\chapter{Granger e Toda-Yamamoto}
\label{cap:granger}

\section{O que o teste de Granger mede, e o que o nome sugere}

O teste de Granger mede \textbf{precedência preditiva}: se o passado de uma
série ajuda a prever a outra, além do que o passado da própria já previa. Ele
\textbf{não} mede causalidade, apesar do nome consagrado de ``causalidade de
Granger''.

A forma da regressão é
\begin{equation}
y_t = c + \sum_{j=1}^{p} a_j\, y_{t-j} + \sum_{j=1}^{p} b_j\, x_{t-j} + u_t,
\end{equation}
e testa-se, por teste F ou Wald, a hipótese conjunta $b_1 = \dots = b_p = 0$.
Rejeitá-la é dizer que o passado de $x$ acrescenta poder preditivo sobre $y$.

\begin{nabanca}
``Uso o nome com rigor: o que ele mede é precedência preditiva. Um teste de
Granger nunca autoriza a palavra `causa' no meu texto.''
\end{nabanca}

\section{Toda-Yamamoto: a correção que torna a inferência válida}

O procedimento resolve o problema da integração com um truque elegante. Ajusta
um modelo vetorial autorregressivo \textbf{em níveis} com $p + d_{\max}$
defasagens --- as $p$ defasagens que se quer testar, mais $d_{\max}$ defasagens
extras, iguais à ordem máxima de integração das séries envolvidas --- e aplica
o teste de Wald \textbf{apenas às $p$ primeiras}.

As defasagens extras existem para absorver a não estacionariedade; ao restringir
o teste às primeiras, a estatística de Wald recupera a distribuição
qui-quadrado usual mesmo sob integração.

\begin{ancora}[Aqui, $d_{\max} = 2$]
Fixado pela linha da pastagem do Norte, que só perde a raiz unitária na segunda
diferença. Não é uma suposição: é o que a tabela de integração mostra.
\end{ancora}

\section{As três colunas de $p$, que não são a mesma régua}

O apêndice reporta três, e a distinção é o próprio argumento:

\begin{itemize}
\item \textbf{$p$ Granger} --- o teste F clássico, em primeira diferença;
\item \textbf{$p$ Wald HAC} --- o mesmo teste com covariância de Newey-West.
Corrige o \emph{erro-padrão}, e \textbf{não} a integração;
\item \textbf{$p$ Toda-Yamamoto} --- o único válido sob integração, e o que o
Capítulo 4 usa para o veredito.
\end{itemize}

\section{O episódio que gerou a decisão D16}

Rodado no sentido inverso, o teste produzia um resultado forte: o Norte
anteciparia o Sul, com $p < 0{,}001$. Se real, ele não salvaria a hipótese do
empurrão --- inverteria a tese inteira.

O diagnóstico veio em três passos, e nenhum deles foi ``a amostra é pequena''.

\vspace{4pt}
\textbf{Primeiro:} a régua estava errada para o material. O teste comparava uma
série já estabilizada com outra que permanece instável mesmo depois de
convertida em variações anuais --- a montagem clássica que fabrica precedência.

\textbf{Segundo:} o método adequado a séries integradas apaga o resultado
\emph{nas duas direções}: $p = 0{,}45$ no sentido Norte--Sul e $p = 0{,}25$ no
sentido Sul--Norte.

\textbf{Terceiro, e mais limpo:} a precedência reaparece onde \emph{não há
mecanismo que a explique}. Em três dos quatro placebos direcionais ela acende, e
o mais eloquente deles é o pasto do Norte ``antecedendo'' o pasto do
\textbf{próprio Sul}, com o mesmo p-valor do resultado original --- o que
nenhuma economia explicaria. O quarto placebo fica em $p = 0{,}056$, e portanto
não acende nem socorre a hipótese.

\begin{quote}\itshape\color{cortinta}
Uma precedência que aparece em qualquer par de séries suaves não é mecanismo, é
defeito de instrumento.
\end{quote}

\begin{nabanca}[O ponto que mais importa neste episódio]
``O veredito é \textbf{sem líder}, e não `o outro lado venceu'. Dizer as duas
direções importa, porque a mesma correção que derruba a precedência do Norte
retira qualquer autorização para eu afirmar que o Sul lidera. A simetria do
veredito é parte do resultado.''
\end{nabanca}

\begin{armadilha}
A tabela do apêndice reporta o bloco reverso \textbf{de propósito}, porque foi
ele que originou a decisão D16, e não porque sustente afirmação. Se a banca
apontar aquele $p < 0{,}001$ na tabela, a resposta é que ele está ali para que
se veja o que o diagnóstico teve de descartar, e que a coluna de Toda-Yamamoto
ao lado é a que decide.
\end{armadilha}

\chapter{Econometria espacial: W, Moran, SLX, SAR e SEM}
\label{cap:espacial}

\section{A matriz de pesos $W$, que é onde tudo começa}

Toda econometria espacial depende de uma definição prévia de ``vizinho'',
codificada numa matriz $W$ de dimensão $n \times n$, com $W_{ij} > 0$ se $j$ é
vizinho de $i$ e $W_{ii} = 0$. Padronizada por linha (cada linha soma 1), o
produto $(W z)_i$ é simplesmente a \textbf{média da variável $z$ nos vizinhos de
$i$}.

O trabalho usa duas construções, e a escolha entre elas é substantiva:

\begin{itemize}
\item \textbf{Contiguidade \emph{queen}} --- vizinhos são as unidades que
compartilham fronteira ou vértice. Usada onde o objeto é a dependência
\emph{genérica}.
\item \textbf{Oito vizinhos mais próximos}, por distância euclidiana entre
centroides em EPSG:5880, com \textbf{filtro direcional}. Em
$W_{\mathrm{sul}}$ só permanecem, entre os oito, os vizinhos de centroide mais
ao \emph{sul}; em $W_{\mathrm{norte}}$, os mais ao \emph{norte}. Usada no teste
do empurrão, onde a direção é a hipótese.
\end{itemize}

\begin{armadilha}[Um detalhe honesto do desenho]
O filtro direcional deixa sem vizinho um punhado de unidades nas bordas: três
linhas inteiramente nulas em $W_{\mathrm{sul}}$ e duas em $W_{\mathrm{norte}}$,
de 166. Elas entram na estimação com termo de vizinhança igual a zero. Está
declarado no apêndice.
\end{armadilha}

\section{O I de Moran: existe organização geográfica no que sobrou?}

\begin{equation}
I = \frac{n}{\sum_i \sum_j W_{ij}} \cdot
\frac{\sum_i \sum_j W_{ij} (z_i - \bar{z})(z_j - \bar{z})}{\sum_i (z_i - \bar{z})^2}.
\end{equation}

Em português: é uma correlação entre o valor de cada unidade e a média dos seus
vizinhos. Positivo significa que vizinhos se parecem; próximo de zero significa
distribuição aleatória no espaço. A significância é obtida por permutação
(embaralha-se aleatoriamente os valores sobre o mapa muitas vezes e compara-se).

Aplicado \textbf{aos resíduos} do painel, ele responde a: ``o que o meu modelo
não explicou ainda tem organização geográfica?''.

\begin{ancora}[A resposta foi sim, e nas duas malhas]
Moran significativo em \textbf{115 de 140} testes na malha municipal e em
\textbf{125 de 140} na malha de AMC. A dependência espacial não é, portanto,
artefato da unidade escolhida.
\end{ancora}

Foi essa constatação que levou a dimensão espacial ao desenho, e que também
obrigou o bootstrap de blocos do Capítulo~\ref{cap:bootstrap}.

\section{As três formas espaciais}

Partindo dos mínimos quadrados com os dois efeitos fixos,
\begin{equation}
\Delta y_{it} = \Delta \mathbf{x}_{it}' \boldsymbol{\beta} + \alpha_i + \lambda_t + \varepsilon_{it},
\end{equation}
cada forma acrescenta um termo em lugar diferente:

\vspace{6pt}
\textbf{SLX --- defasagem espacial dos regressores.} É a do teste do empurrão:
\begin{equation}
\Delta y_{it} = \beta\, \Delta x_{it} + \theta\, (W\Delta x)_{it} + \alpha_i + \lambda_t + \varepsilon_{it}.
\label{eq:slx}
\end{equation}
O parâmetro de interesse é $\theta$, e \textbf{não} $\beta$: $\theta$ mede o
efeito da lavoura \emph{dos vizinhos} sobre o pasto local, ao passo que $\beta$
mede substituição \emph{dentro} da unidade, que não é deslocamento.

\textbf{SAR --- defasagem espacial da variável dependente:}
\begin{equation}
\Delta y_{it} = \rho\, (W\Delta y)_{it} + \Delta \mathbf{x}_{it}' \boldsymbol{\beta} + \alpha_i + \lambda_t + \varepsilon_{it}.
\end{equation}

\textbf{SEM --- dependência no termo de erro:}
\begin{equation}
\Delta y_{it} = \Delta \mathbf{x}_{it}' \boldsymbol{\beta} + \alpha_i + \lambda_t + u_{it},
\qquad u_{it} = \lambda\, (Wu)_{it} + \varepsilon_{it}.
\end{equation}

\subsection{Por que a distinção entre SAR e SEM importa para a leitura}

Sob o \textbf{SAR}, o efeito de um choque numa unidade transborda para as
vizinhas e \emph{volta}, de modo que $\boldsymbol{\beta}$ deixa de ser o efeito
total. Sob o \textbf{SEM}, $\boldsymbol{\beta}$ continua interpretável e o que
muda é a eficiência da estimativa.

Por isso os coeficientes das colunas espaciais \textbf{não são diretamente
comparáveis} ao de mínimos quadrados, e a pergunta que a tabela faz não é se
eles coincidem, e sim se o \emph{sinal} e a \emph{significância} resistem.

\begin{ancora}[O que a modelagem da dependência mudou]
Parâmetros de defasagem e de erro espacial entre $+0{,}35$ e $+0{,}56$
--- a dependência é real e mensurada. \\
Das oito estimativas espaciais, \textbf{sete atenuam-se} em relação aos mínimos
quadrados, no máximo \textbf{14,6\%}, e \textbf{uma cresce 0,6\%}. \\
\textbf{Nenhuma troca de sinal, nenhuma perde significância}: o maior $p$ entre
as oito é da ordem de $10^{-10}$.
\end{ancora}

\begin{nabanca}
``A dependência espacial existe, eu a medi e ela não muda veredito nenhum. O
que a modelagem corrige, no que interessa à leitura, é a barra de erro, e é
por isso que o bootstrap de blocos entrou também nos centros de massa.''
\end{nabanca}

\begin{armadilha}
Não diga ``todas as estimativas espaciais se atenuam''. Sete atenuam; uma cresce
0,6\%. É pequeno e não muda nada, e dizer o número exato é mais forte do que
arredondar para a versão conveniente.
\end{armadilha}

\section{O LISA, e por que ele é exploratório}

Os indicadores locais de associação espacial decompõem o I de Moran unidade a
unidade, mostrando onde estão os aglomerados. Aqui eles são tratados como
\textbf{exploratórios}, por não incorporarem correção para testes múltiplos ---
e são $n$ testes simultâneos, um por unidade.

\section{O resultado do teste do empurrão}

A grade é o produto de três réguas de lavoura (classe Agricultura, união com o
Mosaico, e área de soja do SIDRA) por duas janelas (série plena e truncada em
2019) por três modelos, um deles o placebo com $W_{\mathrm{norte}}$: dezoito
especificações e trinta e seis coeficientes. Das dezoito, as \textbf{doze} não
placebo são as que o Capítulo 4 lê.

\begin{ancora}[O que as doze mostram]
\textbf{Com oito vizinhos, as doze estimativas de $\theta$ são negativas.}
Nenhuma cai na região que a hipótese exige --- ela exigia $\theta > 0$. Com
quatro e com doze vizinhos, algumas do rebanho mudam de lado, sempre com
$p > 0{,}70$: nenhum $\theta$ positivo é significativo nas três matrizes. \\
Quando a lavoura dos vizinhos ao sul cresce, o pasto local \emph{diminui}, em
vez de aumentar. \\
Onze das doze não são significativas; a décima segunda tem $p = 0{,}02$, é
negativa e deixa de ser significativa com quatro ($p = 0{,}29$) e com doze
vizinhos ($p = 0{,}41$).
\end{ancora}

\begin{naodiz}
As doze células \textbf{não trazem cálculo de poder próprio}. O que elas
estabelecem é o \emph{sinal} do coeficiente, e não a exclusão de um empurrão
pequeno. O poder declarado é o do teste temporal. \\[4pt]
E o placebo direcional não é uniformemente nulo: numa das seis células (régua
da soja do IBGE, janela plena) o termo de vizinhança \emph{ao norte} é
significativo a 5\% ($\beta = -0{,}054$; $p = 0{,}032$), onde o termo ao sul
correspondente é nulo. Como o argumento é que um empurrão teria direção, esse
acendimento \textbf{enfraquece a premissa}, e é mais um motivo para não ler
as doze estimativas como um achado.
\end{naodiz}

\section{O que ocupa o lugar do empurrão: a substituição local}

O mesmo desenho revela o maior efeito de toda a frente. É preciso ser exato no
vocabulário: $\theta$ mede o efeito \emph{dos vizinhos}; $\beta$ mede o que
ocorre \emph{dentro do próprio município}. É a diferença entre empurrão, à
distância, e substituição, no mesmo lugar.

\begin{ancora}[A substituição local]
Nas quatro células medidas em área de satélite, o coeficiente vai de
$\mathbf{-0{,}52}$ a $\mathbf{-1{,}14}$, com $p < 0{,}001$ em cinco das seis
células da grade (a exceção é a régua da soja do IBGE na janela curta, com
poucos anos). \\
As duas células que regridem sobre a área de soja declarada ao IBGE ficam em
torno de $-0{,}03$ a $-0{,}07$.
\end{ancora}

\begin{armadilha}[Duas comparações que não se podem fazer]
\textbf{Primeira:} os $-0{,}03$ da soja \emph{não} são comparáveis em magnitude
aos $-0{,}52$ da agricultura, porque o regressor é outro --- um hectare de soja
é subconjunto de um hectare de agricultura. O que se compara entre réguas é o
sinal e a significância, não o tamanho. \\[4pt]
\textbf{Segunda:} o coeficiente dobrar de $-0{,}52$ para $-1{,}14$ na régua da
união \textbf{não é correção}. A união é o teto do intervalo previsto na D26,
não a medida certa. E um coeficiente de $-1{,}14$ não descreve substituição
``mais completa'': descreve mais de um hectare de pasto perdido por hectare de
lavoura ganho, o que por definição \emph{excede} a substituição e inclui algum
fluxo adicional não identificado neste desenho.
\end{armadilha}

\chapter{O desenho \emph{shift-share} e a inferência por permutação}
\label{cap:shiftshare}

\section{A ideia}

A hipótese é que um mesmo motor macroeconômico atinja de modo diferente lugares
com exposição diferente. O desenho de interação do tipo \emph{shift-share} (ou
Bartik) tem dois ingredientes:

\begin{itemize}
\item o \textbf{\emph{shift}} --- o choque nacional, comum a todas as unidades
em cada ano. Aqui, a variação padronizada do índice de taxa de câmbio efetiva
real, defasada um ano;
\item a \textbf{\emph{share}} --- a exposição local, fixa no tempo. Aqui, o
gradiente de aptidão edafoclimática da Embrapa, padronizado em escore $z$.
\end{itemize}

\section{A equação, e as três propriedades que a governam}

\begin{equation}
\Delta y_{it} = \gamma \left( \tilde{s}_{t-1} \times \tilde{E}_i \right)
              + \alpha_i + \lambda_t + \varepsilon_{it}
\label{eq:shiftshare}
\end{equation}

com $\Delta y_{it}$ = variação anual do efetivo bovino da unidade $i$, em
milhões de cabeças.

\vspace{6pt}
\textbf{Propriedade 1: não há termos principais.} O \emph{shifter} é comum a
todas as unidades num dado ano e é absorvido por $\lambda_t$; a exposição é fixa
no tempo e é absorvida por $\alpha_i$. O que resta identificado é apenas o
\textbf{gradiente} do efeito --- se o choque morde mais onde a exposição é maior
--- e \emph{nunca o seu nível}.

\textbf{Propriedade 2:} por isso mesmo, $\gamma$ \textbf{não mede o efeito do
câmbio sobre o rebanho}, e nenhum número do apêndice autoriza essa leitura.

\textbf{Propriedade 3:} o número de realizações \emph{independentes} do choque é
o número de \emph{anos}, cerca de 38, e não o número de células do painel.
Nenhuma quantidade adicional de unidades espaciais o aumenta.

\begin{nabanca}
``Este desenho não estima o efeito do câmbio. Ele estima a geografia desse
efeito dentro do estado. O nível causal do canal cambial vem da literatura, em
escala nacional e continental; a minha contribuição possível é onde ele bate
mais forte.''
\end{nabanca}

\section{A previsão, que é contraintuitiva e vale explicar}

Sob depreciação cambial, a previsão é que o rebanho cresça \textbf{mais onde a
aptidão edafoclimática é baixa} --- a fronteira do Norte --- e menos no núcleo
apto do Sul, onde a mesma terra tem uso alternativo mais rentável em lavoura. O
coeficiente esperado da interação é, portanto, \textbf{negativo}.

\section{A camada de aptidão, e as três propriedades declaradas}

A camada é o mapa de aptidão agrícola das terras do Brasil, da Embrapa,
cartografia vetorial na escala 1:500.000, com 8.284 polígonos recortados em
Goiás. O dígito líder do símbolo é o grupo de aptidão, de 1 (a melhor para
lavouras) a 6 (indicada para preservação), convertido em escore por
$7 - \text{grupo}$, de modo que o valor cresce com a aptidão. A cobertura
efetiva é de 33,0 Mha, ou 97\% do estado.

\begin{ancora}[O gradiente aparece sozinho]
Sul \textbf{4,69}, Centro \textbf{4,47}, Norte \textbf{4,17} na escala ordinal
da fonte, com as três regiões diferindo a $p = 0{,}0002$. \\
Correlação com a latitude: $\mathbf{-0{,}44}$ (Pearson), $-0{,}40$ (Spearman).
\end{ancora}

\textbf{Primeira propriedade --- é o que a torna útil.} Ela é construída de
solo, clima e relevo, e não do uso observado. Isso resolve a
\textbf{circularidade}: a ordenação Sul--Centro--Norte reaparece sem que a
variável a explicar entre na conta.

\textbf{Segunda --- ela é estática.} Um valor por unidade, sem variação no
tempo. Toda a identificação vem da interação com o choque, e não de mudança na
exposição.

\textbf{Terceira --- uma suposição de medida que convém declarar.} O escore
$7 - \text{grupo}$ trata como \emph{cardinal}, com espaçamento igual entre
grupos, o que na origem é uma classificação \emph{ordinal}. Uma construção
alternativa que não faz essa suposição --- a fração da unidade em grupos aptos
à lavoura --- é calculada da mesma camada e acompanha a média do escore na
medida de qualidade do remanescente.

\begin{naodiz}
A escala 1:500.000 é grosseira diante dos 30 m da série de cobertura: a camada
resolve \textbf{gradiente regional}, e não contraste local.
\end{naodiz}

\section{Por que o erro-padrão agrupado é otimista aqui}

Num desenho \emph{shift-share} com um \textbf{único} \emph{shifter}, o
erro-padrão agrupado trata cada ano como uma realização independente do choque.
Mas há um só choque, observado 38 vezes, e ele é serialmente correlacionado. O
resultado é um erro-padrão pequeno demais e um $p$ otimista.

\section{A inferência por permutação, passo a passo}

A decisão D20 substitui a régua. O procedimento:

\begin{enumerate}
\item Mantenha fixas as \emph{shares} $\tilde{E}_i$, o painel e os efeitos
fixos.
\item Reembaralhe a série anual do câmbio entre os anos.
\item Recalcule a interação e o coeficiente $\hat{\gamma}$.
\item Repita, e conte quantas realizações nulas igualam ou superam o observado
em módulo. Essa fração é o $p$.
\end{enumerate}

\textbf{Duas variantes, e a distinção é relevante.} A permutação \emph{livre}
sorteia qualquer ordem e assim \textbf{quebra a autocorrelação} do câmbio,
tendendo a ser otimista. A permutação \textbf{circular} apenas
\emph{rotaciona} a série sobre suas posições, o que \textbf{preserva} a
autocorrelação do choque, e é a versão defensável para uma série serialmente
correlacionada.

\section{O piso de $1/38$, que é o detalhe mais elegante do trabalho}

São usadas \emph{todas} as $T-1$ rotações possíveis, e não uma amostra delas. E
o $p$ reportado \textbf{inclui a estatística observada entre as realizações}:
é o número de rotações ao menos tão extremas quanto a observada, \emph{mais
um}, dividido por $T$.

Essa convenção --- e não a proporção simples, que teria piso zero --- faz o
menor $p$ atingível ser
\begin{equation}
p_{\min} = \frac{1}{38} \approx 0{,}026 .
\end{equation}

\begin{quote}\itshape\color{cortinta}
A resolução do teste, e não apenas o seu resultado, é parte do que se pode
afirmar.
\end{quote}

Por isso, quando a latitude devolve $p = 0{,}026$, isso significa que
\textbf{nenhuma rotação superou o observado} --- o melhor desfecho possível do
teste --- e não uma margem folgada. E o $0{,}053$ da especificação seguinte é
decidido por \emph{uma única rotação}.

\begin{nabanca}
``Esse 0,026 é o piso da minha régua, e não uma folga. Com 38 realizações do
choque, o menor p-valor que a permutação circular pode devolver é 1 sobre 38. O
que ele me diz é que nenhuma rotação superou o observado; ele não me dá margem
para tratar isso como significância robusta.''
\end{nabanca}

\section{O resultado, e as duas camadas de qualificação}

\begin{ancora}[A especificação confirmatória]
$\beta = \mathbf{-0{,}033}$, com $R^2$ \emph{within} de $0{,}001$. \\
$p$ agrupado $= 0{,}026$; $p$ de permutação entre \textbf{0,07 e 0,13}. \\
\textbf{Veredito: corroborante, não estabelecido.}
\end{ancora}

O que sustenta a leitura não é a significância, e sim a \textbf{especificidade}:
cada teste que deveria sair vazio sai vazio.

\begin{itemize}
\item Placebos de desfecho: área urbana $p = 0{,}34$; água $p = 0{,}33$.
\item Placebos no tempo, que adiantam o \emph{shifter} de modo que o efeito
precederia a causa: $p = 0{,}105$ para a liderança de um ano, $p = 0{,}77$ para
a de dois.
\item \emph{Jackknife} por ano: o sinal é estável em \textbf{38 dos 38} anos.
\end{itemize}

\begin{armadilha}[Não cite só os placebos mais vazios]
Os oito placebos não são igualmente vazios. Os dois que chegam mais perto de
rejeitar são os de tempo da \emph{outra} exposição, a de fronteira, com $p$ de
$0{,}062$ e $0{,}070$. Eles não socorrem a hipótese, porque o \textbf{sinal
deles é o oposto} do que ela prevê --- um placebo que quase acende na direção
contrária não é evidência a favor. Mas citar apenas os dois mais nulos seria
escolher a régua depois de ver o resultado.
\end{armadilha}

\begin{armadilha}[E não compare réguas diferentes]
O $p$ do achado é o de \emph{permutação}, o conservador; os dos placebos são os
\emph{agrupados}, que se sabe otimistas. A permutação não foi computada para as
linhas de placebo. Comparar $0{,}105$ com $0{,}07$--$0{,}13$ seria tomar por
igual o que foi medido de dois modos.
\end{armadilha}

\section{A corrida entre exposições, e a decisão D28}

Aqui está a autocorreção mais consequente da frente, e ela precisa ser
apresentada por você antes de a banca chegar nela.

Toda a bateria acima é de placebos de \textbf{desfecho}: ela mostra que o sinal
aparece na variável que deveria responder e não nas outras. Ela não faz a
pergunta decisiva para uma \emph{share} espacial: \emph{a exposição certa é
mesmo a aptidão?}

\begin{quote}\itshape\color{cortinta}
Exógena não é o mesmo que isolada.
\end{quote}

A aptidão correlaciona-se $-0{,}44$ com a latitude, e a latitude organiza, em
Goiás, quase tudo o que varia de sul para norte: idade da ocupação,
infraestrutura, distância aos mercados, especialização produtiva, preço da
terra. Se o choque comum interage com \emph{qualquer} gradiente Sul--Norte, o
coeficiente da aptidão acende sem que a aptidão seja o mecanismo.

O teste é pôr as candidatas na \textbf{mesma regressão}. Além da aptidão,
entram a \textbf{latitude} do centroide (o confundidor puro, que mede ``quão ao
norte'' sem nenhum conteúdo agronômico) e a \textbf{distância ao núcleo de
lavoura de 1985} (aproximação de acesso, construída de dado anterior a toda a
janela estimada).

O fator de inflação de variância entre aptidão e latitude é \textbf{1,24}: a
correlação é moderada e não colinearidade, de modo que as duas são separáveis em
princípio e a corrida é informativa.

\vspace{6pt}
\begin{tabular}{@{}llrrr@{}}
\toprule
\textbf{Esp.} & \textbf{Exposição} & \textbf{$\hat{\gamma}$} & \textbf{$p_{\mathrm{agr}}$} & \textbf{$p_{\mathrm{perm}}$} \\
\midrule
S1 & Aptidão sozinha & $-0{,}033$ & $0{,}026$ & $0{,}132$ \\
S2 & Latitude sozinha & $+0{,}051$ & $0{,}015$ & $\mathbf{0{,}026}$ \\
S3 & Acesso sozinho & $+0{,}028$ & $0{,}118$ & $0{,}184$ \\
\addlinespace
S4 & Aptidão (com latitude) & $-0{,}012$ & $0{,}302$ & $0{,}474$ \\
S4 & Latitude (com aptidão) & $+0{,}046$ & $0{,}028$ & $0{,}053$ \\
\addlinespace
S5 & Aptidão (com acesso) & $-0{,}024$ & $0{,}139$ & $0{,}132$ \\
S5 & Acesso (com aptidão) & $+0{,}015$ & $0{,}482$ & $0{,}395$ \\
\addlinespace
S6 & Aptidão (com as duas) & $-0{,}026$ & $0{,}111$ & $0{,}132$ \\
S6 & Latitude (com as duas) & $+0{,}081$ & $0{,}052$ & $0{,}079$ \\
S6 & Acesso (com as duas) & $-0{,}051$ & $0{,}236$ & $0{,}237$ \\
\bottomrule
\end{tabular}

\vspace{10pt}
\textbf{A leitura.} Sozinha, a latitude entrega um gradiente \emph{mais forte}
que o da aptidão, e é a única das três exposições que cruza o corte de 5\% sob a
régua conservadora. Postas juntas (S4), a aptidão perde \textbf{62\% da
magnitude} ($-0{,}033 \to -0{,}012$) e o seu $p$ de permutação sobe de $0{,}13$
para $0{,}47$; na régua agrupada, sai de $0{,}026$ para $0{,}30$, isto é, deixa
de ser significativa \emph{mesmo sob a inferência otimista}. Enquanto isso a
latitude quase não se move.

\subsection{A especificação S6, que é a que menos favorece a leitura adotada}

Vale conhecê-la bem, porque uma banca atenta vai encontrá-la. Com as três
exposições disputando o mesmo choque, a aptidão \emph{recupera} parte da
magnitude ($-0{,}026$) e a latitude \emph{deixa} de cruzar os 5\%
($p_{\mathrm{perm}} = 0{,}079$).

\begin{nabanca}[A resposta exata para a S6]
``Ela não desfaz a D28, porque a decisão se apoia em a aptidão \textbf{perder}
significância diante do confundidor puro, e não em a latitude \textbf{ganhá-la}.
O que a especificação com as três mostra é a mesma coisa por outro lado: três
gradientes correlacionados e cerca de 38 realizações do choque não se separam
entre si. E é esse o veredito que eu reporto.''
\end{nabanca}

\section{O que a decisão D28 fixa para a redação}

\begin{quote}\itshape\color{cortinta}
Onde o texto disser ``gradiente de aptidão'', leia-se a exposição
\emph{medida} por aptidão --- útil porque é exógena ao uso da terra,
insuficiente para nomear o canal. Nenhuma afirmação do texto atribui o efeito a
solo e clima em particular.
\end{quote}

Os dois sinais dizem a mesma coisa em vocabulários diferentes, porque ao norte
a aptidão é baixa. O que a corrida decide é qual rótulo o dado sustenta, e ele
sustenta o do \textbf{eixo}, não o do solo.

\section{As três qualificações que puxam o peso do achado para baixo}

Diga-as você mesmo, na ordem:

\begin{enumerate}
\item \textbf{Os desfechos de área da mesma grade são nulos.} A interação entre
câmbio e aptidão não move a variação de área de pastagem ($p = 0{,}92$ sem
defasagem, $0{,}83$ com um ano), e a interação entre preço da soja e aptidão não
move a variação de área de agricultura ($p = 0{,}77$ e $0{,}60$). Como a
narrativa é sobre procura por terra, é preciso registrar que o único gradiente
com sinal aparece numa \emph{contagem de cabeças}, e não na terra.
\item \textbf{Multiplicidade.} A célula vem de uma grade de \textbf{192} testes
(quatro exposições, quatro \emph{drivers}, quatro desfechos, três defasagens), e
o exame de sensibilidade mostrou que ela sobrevive ao controle de falsas
descobertas numa família de \textbf{96} testes e \emph{não} sobrevive numa de
\textbf{144}. É sensível ao recorte da família.
\item \textbf{O desenho não poderia estimar o nível.} Com efeito fixo de ano, só
o gradiente é estimável, e é sobre o gradiente, não sobre o câmbio como
causa, que a evidência é corroborante.
\end{enumerate}

\section{O teto estrutural desta frente}

\begin{naodiz}
O poder é limitado pela dimensão \textbf{temporal}, e não aumenta com o número
de unidades espaciais. Levantar o teto exigiria: um \emph{shifter} com variação
transversal própria (exposição diferencial a preços por produto, por exemplo,
em vez de uma única série nacional), \emph{ou} uma janela temporal maior,
\emph{ou} um painel de vários estados que multiplicasse as realizações do
choque. São \textbf{desenhos distintos}, e não robustez adicional sobre este.
\end{naodiz}

\chapter{Fluxo, taxa e estoque: a matemática do teto de oferta}
\label{cap:oferta}

\section{A identidade, e por que ela precisa ser declarada}

O fluxo de conversão é, \emph{por construção}, o produto entre o estoque ainda
exposto e a taxa com que ele é convertido:
\begin{equation}
\text{fluxo}_{it} \equiv \text{taxa}_{it} \times \text{estoque}_{i,t-1},
\qquad \text{com}\quad \text{taxa}_{it} = \frac{\text{fluxo}_{it}}{\text{estoque}_{i,t-1}}.
\end{equation}

Segue daí, \textbf{sem nenhum dado}, que se a taxa for aproximadamente constante
o fluxo acompanha o estoque.

\section{As duas equações, e qual delas tem conteúdo}

\begin{equation}
\text{fluxo}_{it} = \beta\, \text{estoque}_{i,t-1} + \alpha_i + \lambda_t + \varepsilon_{it}
\label{eq:b1}
\end{equation}
\begin{equation}
\text{taxa}_{it} = \beta\, \text{esgotamento}_{i,t-1} + \alpha_i + \lambda_t + \varepsilon_{it}
\label{eq:b2b}
\end{equation}

com $\text{esgotamento}_{it} = 1 - \dfrac{\text{estoque}_{it}}{\text{estoque}_{i,1985}}$.

\vspace{6pt}
A Equação~\ref{eq:b1} carrega a identidade e \textbf{não pode surpreender}. A
regressão devolve exatamente o esperado ($\beta = +2{,}76$, com $R^2$
\emph{within} de $0{,}43$), e esse número \textbf{não é evidência de coisa
alguma}: é a identidade se reapresentando.

Acrescentar um termo quadrático no estoque não muda nada --- ele sai nulo
($p = 0{,}92$) ---, o que diz que a relação é a linha reta que a identidade
prevê, sem curvatura própria. É uma afirmação sobre o \emph{estoque}, e não
sobre o comportamento da taxa.

\begin{nabanca}[Dizer isto antes de a banca dizer]
``A primeira regressão não é um resultado meu. Ela é uma identidade contábil se
reapresentando, e eu a reporto justamente para não deixar que ela passe por
achado. O conteúdo empírico está na segunda.''
\end{nabanca}

\section{O conteúdo empírico: a taxa cai com o esgotamento}

Aqui as equações \emph{não} entram em primeira diferença, e essa é uma das
três exceções à D7. A razão é que o que se mede é o que acontece a uma unidade
\textbf{enquanto ela se esgota}, e diferenciar apagaria justamente esse nível. A
tendência é retirada por efeitos fixos de unidade e de ano.

\begin{ancora}[O resultado, e o que o sustenta]
Numa grade de \textbf{dezesseis} combinações (quatro tratamentos da variável de
esgotamento, com e sem ponderação pelo tamanho do estoque, em duas amostras), o
coeficiente é \textbf{negativo nas dezesseis} e cruza o corte de 5\% em
\textbf{onze}. \\[4pt]
Nas combinações que põem a variável dentro do domínio, os tratamentos convergem
para \textbf{0,5 a 0,8 ponto percentual de taxa anual a menos a cada dez pontos
de esgotamento}. \\[4pt]
A fração da variação explicada dentro das unidades, nula na especificação
original, passa a situar-se entre \textbf{4\% e 20\%}.
\end{ancora}

\section{A decisão D29, e por que o nulo anterior era artefato}

O nulo publicado antes dizia que a taxa \emph{não} cairia com o esgotamento.
Ele era artefato, e a causa é instrutiva.

A variável de esgotamento é uma fração do estoque de 1985 e deveria viver entre
zero e um. No painel ela alcançava $\mathbf{-84{,}9}$, em \textbf{14\%} dos
pares unidade-ano --- vindos de unidades cujo estoque inicial era próximo de
zero e cujo estoque \emph{cresceu} ao longo da série por oscilação de
classificador.

Padronizada, a variável passou a ser dominada por esses poucos valores, e o
coeficiente foi achatado contra zero. A razão é aritmética: o escore $z$ divide
pelo desvio-padrão, e um único valor oitenta vezes fora de escala redefine a
unidade de medida de toda a coluna.

\begin{quote}\itshape\color{cortinta}
Um z-score não é robusto a nada.
\end{quote}

\subsection{As três regras que a D29 fixa, e que valem além deste caso}

\begin{enumerate}
\item \textbf{Toda variável com domínio declarado é verificada contra ele antes
de entrar padronizada.}
\item \textbf{Onde há tratamento possível, o resultado reportado é o que
concorda entre os tratamentos}, e não o de um deles. A correção foi aceita
porque as dezesseis combinações devolvem o mesmo sinal, e não porque a régua
nova fosse mais conveniente que a antiga.
\item \textbf{A verificação de um resultado herda a régua de inferência do teste
verificado.} Trocá-la por uma régua mais frouxa no meio da conferência
produziria uma robustez que é da régua, e não do dado.
\end{enumerate}

\begin{armadilha}[A comparação entre tratamentos exige unidade natural]
O desvio-padrão da variável tratada varia \textbf{dezessete vezes} de um
tratamento a outro (de $0{,}21$ no domínio a $3{,}54$ sem tratamento). Portanto
os coeficientes \emph{padronizados} \textbf{não são comparáveis entre si}. É o
coeficiente em unidade natural --- pontos percentuais de taxa por décimo de
esgotamento --- que autoriza dizer que concordam.
\end{armadilha}

\section{As três ressalvas que a concordância não dissolve}

\begin{enumerate}
\item O tratamento por \textbf{domínio}, o mais limpo conceitualmente, descarta
\textbf{46 unidades}, e elas detêm \textbf{17,3\% do estoque convertível de
2024}. Para a regressão a exclusão é correta, porque nessas unidades o
esgotamento é \emph{indefinido}, e não extremo; para quem for somar hectares,
não é inócua.
\item O tratamento alternativo, o \textbf{piso em zero}, evita o descarte ao
custo de criar um acúmulo artificial de observações exatamente em zero. Os dois
têm defeitos \emph{diferentes} e chegam ao mesmo lugar, e é essa a razão de a
leitura descansar na concordância.
\item Nenhum dos quatro conserta a \emph{causa} de as 46 existirem, que é a
oscilação do classificador entre pastagem e savana.
\end{enumerate}

\section{Multiplicidade: por que as dezesseis não recebem correção}

\begin{nabanca}
``As dezesseis células são dezesseis leituras da \emph{mesma} hipótese,
pré-declarada antes da estimação, e não uma família de hipóteses distintas. O
controle de falsas descobertas incide sobre grades exploratórias, que perguntam
coisas diferentes; aplicá-lo a uma grade de robustez, que pergunta a mesma coisa
de quatro maneiras, puniria a checagem em vez da pescaria.''
\end{nabanca}

\section{A decomposição entre estoque e taxa, e a leitura que a inverte}

A mudança do fluxo entre dois atos separa-se aritmeticamente em duas parcelas:
quanto vem de \emph{haver menos Cerrado a converter} e quanto vem de
\emph{converter a uma taxa diferente} o Cerrado que ainda há.

\begin{ancora}[O resultado, e a armadilha embutida]
A parcela de estoque é negativa nas três regiões. \\
No Sul, apenas \textbf{17\%} da freada é, mecanicamente, ter menos Cerrado. Os
outros \textbf{83\%} estão na \emph{parcela residual}.
\end{ancora}

\begin{armadilha}[A leitura descuidada diz o contrário do que a evidência sustenta]
A parcela residual \textbf{não é demanda}. Ela reúne tudo o que não é o tamanho
do estoque: propensão a converter, atrito de acesso, proteção, e a troca de
fonte de terra. Se o argumento de oferta dependesse do tamanho da parcela de
estoque, ele seria fraco, e ele não depende.
\end{armadilha}

\section{Por que o argumento se monta por eliminação, e não pela decomposição}

O argumento tem três pernas, e nenhuma delas é o tamanho da parcela de estoque:

\begin{enumerate}
\item \textbf{Não é falta de comprador.} A taxa do Sul caiu enquanto a demanda
estava no pico.
\item \textbf{Não é a lei.} A proteção integral não estava no caminho (ver
adiante).
\item \textbf{A demanda de terra do Sul continuou existindo e foi atendida por
pasto já aberto.}
\end{enumerate}

A isso se soma a regularidade medida --- a taxa caindo com o esgotamento --- e
uma segunda medida independente, a de que o remanescente do Sul é hoje, em três
quintos, a fisionomia que a fronteira não consome.

\section{O sinal do viés, que corre a favor da leitura}

Detalhe fino e forte. Taxa e esgotamento compartilham o estoque anterior --- no
denominador de uma, no numerador com sinal trocado da outra --- e esse
acoplamento empurra o coeficiente para \textbf{cima}.

\begin{quote}\itshape\color{cortinta}
Achar sinal negativo é achar \emph{contra} a montagem, e não a favor dela.
\end{quote}

\begin{naodiz}
O que a grade mede é que, quando uma unidade se esgota, \textbf{caem as duas
coisas ao mesmo tempo}: o estoque que resta e a taxa com que ele é convertido.
Isso é uma \emph{regularidade compatível com atrito de oferta}, e não a
demonstração de que a dificuldade do remanescente seja a \emph{causa} da queda
--- esgotamento e taxa são construídas a partir do mesmo estoque, e os efeitos
fixos reduzem o confundimento sem transformar uma relação dinâmica em mecanismo
identificado.
\end{naodiz}

\chapter{Multiplicidade e poder}
\label{cap:multiplicidade}

\section{O problema, em uma frase}

Com corte de 5\%, cerca de cinco em cada cem testes aparecem significativos por
acaso. Algumas grades deste trabalho passam de uma centena de testes.

\section{Benjamini-Hochberg: controle da taxa de falsas descobertas}

O procedimento ordena os $m$ p-valores em ordem crescente,
$p_{(1)} \leq \dots \leq p_{(m)}$, e encontra o maior $k$ tal que
\begin{equation}
p_{(k)} \leq \frac{k}{m}\, q ,
\end{equation}
rejeitando as $k$ hipóteses correspondentes. O parâmetro $q$ é a fração de
falsos positivos que se aceita \emph{entre as descobertas}.

A diferença em relação a Bonferroni vale ser dita: Bonferroni controla a
probabilidade de \emph{qualquer} falso positivo e é muito conservador; o FDR
controla a \emph{proporção} de falsos entre os rejeitados, e é o adequado para
grades exploratórias.

\begin{quote}\itshape\color{cortinta}
O tamanho da família altera o veredito. Um resultado que sobrevive numa grade
pequena pode não sobreviver quando ela é ampliada, e por isso o trabalho
declara explicitamente quantos testes compõem a família em cada caso.
\end{quote}

É exatamente o que acontece com o \emph{drive} comum: sobrevive numa família de
96, não sobrevive numa de 144.

\section{Poder estatístico, e por que declará-lo é obrigatório num nulo}

O poder é a probabilidade de detectar um efeito \emph{que existe}, dado o
tamanho da amostra e o tamanho do efeito. Um teste com poder baixo que não
rejeita não informou quase nada.

\begin{ancora}[Os poderes declarados neste trabalho]
Teste temporal do empurrão, com 38 diferenças anuais: descarta um empurrão
\textbf{forte} com poder aproximado de \textbf{93\%}; não descarta um
\textbf{moderado}, com poder aproximado de \textbf{48\%}. \\[4pt]
Quarto período candidato (2001--2005 contra 2006--2019): $p = 0{,}060$, com
poder de \textbf{0,63} para as quatro observações disponíveis.
\end{ancora}

\begin{nabanca}
``Sozinho, o teste temporal não fecharia a questão --- ele tem poder de 48\%
contra um empurrão moderado. Somado às doze especificações espaciais, que dão
todas o sinal contrário ao exigido, ele completa o quadro. É a convergência de
duas assinaturas independentes que sustenta o nulo, e não o p-valor de nenhuma
delas.''
\end{nabanca}

\chapter{A conta de carbono}
\label{cap:carbono}

\section{A aritmética}

A conta inteira é área medida multiplicada por densidade de tabela:
\begin{equation}
\text{estoque removido} = \sum_{f} A_f \times \rho_f ,
\end{equation}
onde $A_f$ é a área convertida da formação $f$ e $\rho_f$ a densidade de carbono
daquela formação. A conversão para CO\textsubscript{2} equivalente usa o fator
estequiométrico $44/12$ (a razão entre as massas moleculares do CO$_2$ e do
carbono).

\section{As duas grandezas, que a decisão D31 obriga a distinguir}

\textbf{Estoque removido} é a soma acima: o carbono que estava na cobertura de
origem.

\textbf{Emissão líquida} é o método de diferença de estoque do IPCC em seu nível
mais simples: a diferença entre o estoque \emph{antes} e o estoque
\emph{depois}, porque a pastagem e a lavoura que entram no lugar não estocam
carbono nenhum apenas em teoria.
\begin{equation}
\text{emissão líquida} = \sum_f A_f \times (\rho_f - \rho_{\text{destino}}).
\end{equation}

\begin{ancora}[Os números]
Estoque removido: \textbf{973 Mt} de CO\textsubscript{2}e. \\
Emissão líquida: \textbf{833 Mt} --- o desconto retira \textbf{14,4\%}. \\
Densidade de destino: \textbf{6,63 tC/ha}, vinda da composição observada do
fluxo (64,1\% pastagem, 32,9\% Mosaico, 2,7\% agricultura, 0,3\% urbana). \\
Convenção alternativa (todo o Mosaico como pastagem): densidade 7,48, desconto
16,2\%, total 815 Mt --- \emph{sensibilidade a uma convenção declarada}, e não
incerteza da fonte.
\end{ancora}

\section{As densidades, e por que a procedência importa (decisão D30)}

Os estoques vêm do \textbf{Quarto Inventário Nacional}, que é a régua com que o
país reporta à Convenção do Clima. Três propriedades justificam a escolha: são
publicados por \emph{fitofisionomia}, e não por categoria genérica; no Cerrado
são \emph{diferenciados por estado}, e Goiás tem valor próprio; e a agregação
nas classes de cobertura usadas aqui já está feita e publicada, de modo que
nenhuma correspondência precisou ser arbitrada.

\vspace{6pt}
\begin{tabular}{@{}lrr@{}}
\toprule
\textbf{Formação} & \textbf{tC/ha} & \textbf{Área perdida (Mha)} \\
\midrule
Formação florestal (galeria e cerradão) & 64,72 & 1,43 \\
Formação savânica (Cerrado \emph{sensu stricto}) & 41,32 & 3,79 \\
Formação campestre & 24,94 & 0,33 \\
Campo alagado e área pantanosa & 36,21 & 0,22 \\
\bottomrule
\end{tabular}

\vspace{10pt}
\subsection{A lição de método que veio junto, e que vale para toda a dissertação}

Este é um dos pontos mais elegantes do trabalho, e ele generaliza.

Sob a régua anterior, a afirmação era que a formação florestal emitia mais que a
savânica, apesar de perder menos área. Essa afirmação \textbf{não dependia das
seis densidades}: dependia de \emph{uma razão} entre duas delas superar
\textbf{2,64} --- a razão entre as áreas perdidas. A régua antiga entregava
2,88, com nove por cento de folga. E os três cenários de sensibilidade moviam as
duas densidades \emph{juntas}, mantendo a razão entre 2,88 e 3,00.

\begin{quote}\itshape\color{cortinta}
Eles testavam o nível e eram, por construção, incapazes de testar a composição.
\end{quote}

Sob a régua oficial a razão é \textbf{1,57}, e a afirmação cai.

\begin{nabanca}[A regra geral]
``Quando uma conclusão depende de uma razão entre parâmetros, é a \emph{razão}
que precisa entrar na análise de sensibilidade, e não cada parâmetro
isoladamente. Uma sensibilidade que move os parâmetros na mesma direção não é
sensibilidade nenhuma para a afirmação que depende do contraste entre eles.''
\end{nabanca}

\section{A corroboração externa que a troca revelou}

Vale notar que o \textbf{total praticamente não se moveu}: 973 Mt sob as duas
réguas. Isso é corroboração externa da magnitude por uma fonte independente. O
que se inverteu foi a \emph{composição}.

\section{O escopo, e os dois limites que vão em direções opostas}

\textbf{Mais completo que o usual:} o Inventário soma quatro compartimentos ---
biomassa acima do solo, biomassa subterrânea, madeira morta e serrapilheira. No
Cerrado a razão raiz/parte aérea é alta, e a fração subterrânea responde por
quase três quintos do estoque da savana.

\textbf{Menos completo:} o \textbf{carbono do solo continua de fora}, e no
Cerrado ele é parcela grande do total.

\begin{naodiz}
Contada só a biomassa, a conta cairia para algo entre \textbf{82 e 92\%} do
valor, conforme a formação. É uma faixa de \textbf{escopo, não de incerteza}:
ela diz o que foi contado, e não com que precisão. \\[4pt]
Fora da conta ficam também a dinâmica temporal da liberação, os gases que não
são CO\textsubscript{2}, a queima, os produtos madeireiros e a heterogeneidade
de manejo. \textbf{Nenhuma dessas omissões tem sinal garantido}: no Cerrado o
solo sob pastagem bem conduzida pode perder, manter ou ganhar carbono em relação
à vegetação que havia ali. \\[4pt]
Por isso o número \textbf{não é declarado piso nem teto} do passivo. É o que foi
contado, nos compartimentos nomeados.
\end{naodiz}

\section{As duas aproximações declaradas}

\begin{enumerate}
\item A composição do destino é \textbf{estadual e média}: ela não pareia cada
hectare de floresta, savana ou campo ao destino \emph{daquele} hectare, e sim
aplica a mesma densidade média às três formações e aos três atos. Fazer o
pareamento exigiria a matriz de transição resolvida por formação de origem, que
o cubo de sete grupos não separa.
\item O desconto incide sobre o saldo \textbf{líquido} entre as pontas da série,
e não sobre cada conversão datada.
\end{enumerate}

\section{As três verificações que acompanham o desconto}

Porque um desconto quase uniforme por hectare poderia deslocar comparações:

\begin{itemize}
\item a repartição entre os atos \textbf{não se move} (o Ato~I passa de 76,4\%
para 77,0\%);
\item o ordenamento entre formações \textbf{não se inverte} (a razão savânica
sobre florestal vai de 1,69 para 1,58);
\item a comparação por hectare, que é a afirmação substantiva,
\textbf{fortalece-se}: descontar um mesmo valor de dois estoques desiguais
aumenta a razão entre eles, de 1,57 para 1,69.
\end{itemize}
